% Abstract — Positional Role Grammar and Structural Dictionary
% Reconstruction for the Voynich Manuscript
% Approximately 400 words.

The \vms\ (Beinecke MS~408, radiocarbon-dated 1404--1438~CE) has been
shown in three companion studies to resist all tested substitution-based
decipherment models~\citep{jimenez2026negative}, to exhibit a two-tier
positional encoding consistent with an encyclopaedic categorial
code~\citep{jimenez2026twotier}, and to belong to the medieval
astrological-talismanic tradition codified by \textit{Picatrix
Latinus}~\citep{jimenez2026venusring}. Those studies characterised
\emph{what kind of system} the manuscript is and \emph{which tradition}
it belongs to. The present, fourth study asks a different question:
\emph{can the internal combinatorial rules of the code be recovered
without access to the lost key book?}

We report three cumulative findings based on the full 37,087-word
\sta\ 2.0 corpus (166 glyph types, 157,254 tokens, 224 folios).

\textbf{First}, visual glyph families---groups of glyphs sharing a
common graphic root in the \sta\ classification---concentrate in
specific word positions far more than chance predicts. A
frequency-stratified permutation test ($10{,}000$ iterations) yields
$z = -3.38$, $p = 0.0007$, establishing that graphic form encodes
functional role: glyphs that \emph{look alike} tend to \emph{behave
alike} within the word.

\textbf{Second}, this form--function correlation resolves into a
six-role combinatorial grammar. Each of the 34 core glyphs (covering
99.5\% of tokens) is classified as OPENER, ENCODER, CONNECTOR, BODY,
SCAFFOLD, or CLOSER based on its positional profile. The grammar is
strongly directional: OPENER$\to$ENCODER accounts for 63.5\% of
first-to-second-position transitions; several role pairs are effectively
forbidden ($O/E < 0.3$). A leave-one-out folio classifier trained on
role proportions alone achieves 66.4\% accuracy across five manuscript
sections (herbal, biological, astronomical, zodiac, pharmaceutical),
against a 27\% chance baseline. The grammar is \emph{universal} across
sections; what varies is the \emph{proportion} of glyphs per role,
yielding section-specific markers (e.g., glyph D1 is $1.7\times$
enriched in biological folios; Z1 is $2.8\times$ enriched in zodiac
folios).

\textbf{Third}, the grammar induces a five-slot structural schema for
the Voynich word: DOMINIUM (domain classifier, position~1),
GENUS (primary category encoder, position~2), SPECIES (secondary
modifier, position~3, highest entropy at 4.09~bits), CORPUS (body
content, position~4), and FINIS (terminator, positions~5--6). Each slot
has a closed inventory with measurable section specificity (Jensen--Shannon
divergence ranging from 0.186 to 0.226). The schema is exported as a
machine-readable structural dictionary
(\filename{voynich_structural_dictionary.json}).

The reconstructed schema represents the \emph{syntax} of the lost key
book---the combinatorial frame within which the original author assigned
semantic values to each slot. The analysis recovers the rules that
govern which glyph sequences are valid Voynich words without recovering
what those words mean. Content access remains contingent on the recovery
of the key book itself; the structural dictionary delimits the
search space that any future key-book candidate must satisfy.
